\documentclass[a4paper,11pt]{article}
\usepackage[utf8]{inputenc}
\usepackage[vietnamese]{babel}
\usepackage{amsmath,amssymb}
\usepackage{mathptmx}
\usepackage[margin=2.2cm]{geometry}
\usepackage{booktabs}
\usepackage{tabularx}
\usepackage{ragged2e}
\usepackage{enumitem}
\usepackage{microtype}
\usepackage{listings}
\usepackage{xcolor}
\usepackage[hidelinks]{hyperref}

\newcolumntype{Y}{>{\RaggedRight\arraybackslash}X}
\newcolumntype{L}[1]{>{\RaggedRight\arraybackslash}p{#1}}

\lstset{
  basicstyle=\ttfamily\small,
  breaklines=true,
  frame=single,
  backgroundcolor=\color{gray!5},
  rulecolor=\color{gray!40},
  keywordstyle=\bfseries\color{blue!70!black},
  commentstyle=\itshape\color{green!50!black},
  stringstyle=\color{red!60!black},
  showstringspaces=false
}

\title{\textbf{Chương 5: Kiến Trúc Logic \& Schema-First}\\[0.3em]
\large\textit{Chống Trôi Ngữ Cảnh, Đóng Băng Hợp Đồng Dữ Liệu và Máy Trạng Thái Tường Minh}}
\author{\textbf{Giáo Trình Kỹ Nghệ Vibecoding Chuyên Sâu 2026}}
\date{Tháng 10 Năm 2026}

\begin{document}

\maketitle

\begin{abstract}
Hiện tượng ``trôi ngữ cảnh'' (\textit{Context Drift}) và ảo giác kiến trúc là nguyên nhân chủ yếu khiến các dự án vibecoding bị sụp đổ sau vài tuần phát triển: AI quên cấu trúc dữ liệu ban đầu, tự ý đổi tên trường hoặc sinh mã xử lý bất nhất giữa máy khách và máy chủ. Chương 5 xác lập phương pháp luận \textit{Schema-First} để khóa chặt mô hình dữ liệu bằng TypeScript và Pydantic/Zod; chuẩn hóa mô hình máy trạng thái 4 bước (\textit{Explicit 4-State Machine}); thiết lập quy tắc giới hạn mô-đun siêu nhỏ (\textit{Micro-Modules under 200 lines}) và cơ chế ngắt mạch tự động (\textit{Circuit Breaker}) khi tương tác với các dịch vụ AI.
\end{abstract}

\vspace{1em}
\rule{\textwidth}{0.4pt}

\section{Phương pháp luận Schema-First: Khóa chặt hợp đồng dữ liệu}

Trong kỹ nghệ phần mềm truyền thống, lập trình viên thường bắt đầu bằng cách viết logic xử lý hoặc giao diện rồi mới chuẩn hóa kiểu dữ liệu. Khi vibecoding, cách tiếp cận này là một công thức dẫn đến thảm họa:
\begin{itemize}[leftmargin=1.5em]
    \item Nếu không có một hợp đồng dữ liệu bất biến (\textit{Data Contract}), mô hình AI sẽ tự do ``sáng tác'' các cấu trúc JSON khác nhau ở mỗi lượt phản hồi.
    \item Ví dụ: Một hàm backend trả về \texttt{\{ user\_id: 123 \}}, nhưng ở frontend, AI lại viết component chờ đón \texttt{\{ userId: 123 \}} hoặc \texttt{\{ id: 123 \}}. Ứng dụng không báo lỗi cú pháp nhưng giao diện bị trắng xoá hoàn toàn (\textit{Silent Failure}).
\end{itemize}

\subsection{Nguyên tắc Khởi tạo Schema Trước Mã Nguồn}
Quy tắc bất di bất dịch của Mindful Vibecoding: \textbf{Không bao giờ cho phép AI sinh bất kỳ dòng code UI hay API nào trước khi tệp Schema được phê duyệt và lưu vào đĩa}.

\begin{lstlisting}[language=bash, caption={Ví dụ hợp đồng dữ liệu bất biến bằng Zod và TypeScript}]
// types/quiz.contract.ts
import { z } from 'zod';

export const QuizQuestionSchema = z.object({
  id: z.string().uuid(),
  stem: z.string().min(5),
  options: z.array(z.string()).length(4),
  correctIndex: z.number().int().min(0).max(3),
  explanation: z.string().min(10),
  difficultyLevel: z.enum(['A1', 'A2', 'B1', 'B2', 'C1'])
});

export type QuizQuestion = z.infer<typeof QuizQuestionSchema>;
\end{lstlisting}

\subsection{Người gác cổng tại biên hệ thống (Boundary Gatekeepers)}
Mọi dữ liệu đi vào ứng dụng từ thế giới bên ngoài (yêu cầu HTTP của người dùng, phản hồi từ bên thứ ba, hoặc chuỗi JSON do LLM tạo ra) đều \textbf{bắt buộc phải đi qua cổng xác thực schema (\texttt{safeParse})}:
\begin{enumerate}[label={\arabic*.}]
    \item Nếu dữ liệu hợp lệ: Đưa vào lõi xử lý với kiểu dữ liệu an toàn tuyệt đối.
    \item Nếu dữ liệu bị lỗi hoặc AI sinh thiếu trường: Lập tức chặn đứng và ném lỗi có cấu trúc, không để dữ liệu bẩn lọt sâu vào cơ sở dữ liệu.
\end{enumerate}

\section{Máy trạng thái tường minh 4 bước (Explicit 4-State Machine)}

Một anti-pattern cực kỳ phổ biến trong các ứng dụng AI là việc quản lý trạng thái bằng các cờ boolean rời rạc:
\begin{lstlisting}[language=bash]
// ANTI-PATTERN NGUY HIEM
const [isLoading, setIsLoading] = useState(false);
const [isError, setIsError] = useState(false);
const [data, setData] = useState(null);
\end{lstlisting}

Cách làm này tạo ra các trạng thái bất khả thi (\textit{impossible states}), chẳng hạn như vừa \texttt{isLoading = true} vừa \texttt{isError = true}, khiến giao diện hiển thị đồng thời cả vòng xoay tải trang lẫn thông báo lỗi.

\subsection{Chuẩn hóa 4 trạng thái tương tác hoàn chỉnh}
Mọi component hoặc module kết nối mạng bắt buộc phải được mô hình hóa theo một Máy trạng thái hữu hạn (\textit{Finite State Machine}) gồm 4 trạng thái tách biệt:

\begin{table}[htbp]
\centering
\small
\caption{Chuẩn hóa 4 trạng thái tương tác cho mọi module dữ liệu}
\vspace{0.5em}
\begin{tabularx}{\textwidth}{L{2.5cm} L{4.5cm} L{4.5cm} Y}
\toprule
\textbf{Trạng thái} & \textbf{Điều kiện kích hoạt} & \textbf{Trách nhiệm giao diện} & \textbf{Hành vi kỹ thuật} \\
\midrule
\textbf{1. IDLE} & Mới khởi tạo, người dùng chưa thực hiện hành động & Hiển thị hướng dẫn khởi đầu hoặc khung nhập liệu rỗng & Chưa kích hoạt bất kỳ cuộc gọi mạng nào \\
\textbf{2. LOADING} & Người dùng nhấn gửi yêu cầu, đang chờ mạng hoặc AI xử lý & Hiển thị Skeleton Loader đúng kích thước, khóa nút bấm & Hủy yêu cầu cũ nếu có (\textit{AbortController}) \\
\textbf{3. SUCCESS} & Máy chủ trả về dữ liệu chuẩn khớp $100\%$ với Schema & Render danh sách kết quả, hiển thị dữ liệu chuyển đổi & Nếu danh sách có 0 phần tử, chuyển sang hiển thị \textit{Empty State} \\
\textbf{4. ERROR} & Lỗi mạng, 4xx/5xx hoặc schema validation thất bại & Thông báo nguyên nhân lỗi rõ ràng kèm nút ``Thử lại'' (\textit{Retry}) & Log lỗi nội bộ ra công cụ giám sát, không làm sập toàn bộ app \\
\bottomrule
\end{tabularx}
\end{table}

\begin{lstlisting}[language=bash, caption={Dac ta kieu du lieu cho May trang thai 4 buoc}]
export type AsyncState<T> =
  | { status: 'IDLE' }
  | { status: 'LOADING' }
  | { status: 'SUCCESS'; data: T }
  | { status: 'ERROR'; errorMessage: string; errorCode?: string };
\end{lstlisting}

\section{Định luật Mô-đun siêu nhỏ: Giới hạn dưới 200 dòng mã}

Nghiên cứu về cơ chế chú ý của Transformer chỉ ra rằng: \textbf{Xác suất mô hình AI gây ra lỗi hồi quy (regression bug) hoặc vô tình xóa mất code cũ tăng vọt khi kích thước tệp vượt quá 300 dòng mã}.

Để duy trì độ chính xác của AI ở mức tối đa ($> 95\%$), toàn bộ cơ sở mã phải được tổ chức thành các mô-đun siêu nhỏ (\textit{Micro-Modules}) với nguyên tắc:
\begin{enumerate}[label={\arabic*.}]
    \item \textbf{Ngưỡng cứng 200 dòng}: Bất kỳ tệp mã nguồn nào chạm ngưỡng 200 dòng đều phải được phân rã thành các tệp chuyên biệt:
    \begin{itemize}
        \item \texttt{[feature].types.ts}: Chứa Schema và Interfaces (dưới 50 dòng).
        \item \texttt{[feature].service.ts}: Chứa logic xử lý thuần túy và gọi API (dưới 120 dòng).
        \item \texttt{[feature].ui.tsx}: Chỉ chịu trách nhiệm hiển thị giao diện và nhận props (dưới 150 dòng).
        \item \texttt{[feature].test.ts}: Chứa các ca kiểm thử tự động.
    \end{itemize}
    \item \textbf{Đơn nhiệm tuyệt đối (Single Responsibility)}: Mỗi tệp chỉ giải quyết đúng một việc; khi cần sửa logic dữ liệu, AI chỉ đọc tệp service mà không cần nạp toàn bộ mã nguồn giao diện vào ngữ cảnh.
\end{enumerate}

\section{Cơ chế Ngắt Mạch (Circuit Breaker) cho dịch vụ AI}

Khi ứng dụng gọi các mô hình LLM bên ngoài (Gemini, OpenAI, Claude), các sự cố mạng như nghẽn băng thông, hết hạn mức (\textit{HTTP 429 Rate Limit}), hoặc sập máy chủ (\textit{HTTP 503}) là điều chắc chắn xảy ra.

Một hệ thống Vibecoding chuyên nghiệp bắt buộc phải cài đặt mô hình \textbf{Circuit Breaker}:
\begin{itemize}[leftmargin=1.5em]
    \item \textbf{Trạng thái Closed (Đóng)}: Các yêu cầu được gửi đi bình thường.
    \item \textbf{Trạng thái Open (Mở ngắt mạch)}: Nếu có quá 3 yêu cầu liên tiếp thất bại trong vòng 30 giây, bộ ngắt mạch sẽ tự động bật, tạm dừng mọi yêu cầu mới đến mô hình AI trong 60 giây và trả ngay giao diện dự phòng (\textit{Fallback Cache}) cho người dùng mà không cần chờ timeout.
    \item \textbf{Trạng thái Half-Open (Nửa mở)}: Sau thời gian chờ, gửi một yêu cầu thăm dò; nếu thành công thì đóng mạch trở lại bình thường.
\end{itemize}

\subsection{Cây quyết định: Lựa chọn mô hình kiến trúc logic phần mềm}

Kiến trúc logic không phải là một công thức cố định cho mọi quy mô. Lập trình viên cần đối chiếu bài toán thực tế với 3 mô hình kiến trúc chuẩn mực sau:

\begin{table}[htbp]
\centering
\small
\caption{Cây quyết định lựa chọn mô hình kiến trúc logic phần mềm}
\vspace{0.5em}
\begin{tabularx}{\textwidth}{L{3.2cm} L{3.8cm} L{4.5cm} Y}
\toprule
\textbf{Mô hình kiến trúc} & \textbf{Nguyên lý tổ chức} & \textbf{Công nghệ nòng cốt} & \textbf{Quy mô \& Bối cảnh tối ưu} \\
\midrule
\textbf{Mô hình 1: Micro-Modules Monolith} & Đơn khối phân rã hạt nhân, mỗi tệp $< 200$ dòng, tách biệt schema / service / UI & Next.js / Vite SPA + Server Actions + Zod + SQLite/Supabase & MVP, ứng dụng đơn mục tiêu, quy mô dưới 50.000 người dùng \\
\textbf{Mô hình 2: Serverless Edge Microservices} & Luồng bất đồng bộ, truyền phát trực tiếp (SSE Streaming), ngắt mạch Circuit Breaker & Cloudflare Workers / Vercel Edge + Upstash Redis Queue & Chatbot AI, xử lý hàng đợi, ứng dụng có lưu lượng truy cập đột biến \\
\textbf{Mô hình 3: Local-First Event-Driven} & Xử lý cục bộ 100\%, đồng bộ ngầm khi có mạng, bảo mật tuyệt đối & Tauri / Electron + SQLite FTS5 / PGlite + CRDTs sync & Ứng dụng năng suất cá nhân, ghi chú, xử lý tài liệu nội bộ \\
\bottomrule
\end{tabularx}
\end{table}

\subsection{Kịch bản thực hành đồng hành: state\_machine\_contract\_demo.ts}
Thư mục chương 5 chứa mã nguồn thực hành hoàn chỉnh \texttt{state\_machine\_contract\_demo.ts}. Tệp mã này thể hiện mẫu thiết kế Schema-First với Zod, cài đặt lớp \texttt{AICircuitBreaker} với 3 trạng thái (\textit{Closed, Open, Half-Open}), và cấu trúc Máy trạng thái 4 bước (\textit{IDLE, LOADING, SUCCESS, ERROR}) loại trừ hoàn toàn các trạng thái bất khả thi trong tương tác người dùng.

\section{Tổng kết chương và bài tập thực hành}

\subsection{Tóm tắt cốt lõi}
\begin{itemize}[leftmargin=1.5em]
    \item Tiếp cận Schema-First là cứu cánh duy nhất ngăn chặn hiện tượng trôi ngữ cảnh và bất nhất dữ liệu trong các dự án Vibecoding.
    \item Triệt tiêu toàn bộ boolean \texttt{isLoading/isError} rời rạc; thay thế bằng Máy trạng thái 4 bước tường minh.
    \item Tuân thủ nghiêm ngặt định luật mô-đun siêu nhỏ dưới 200 dòng để bảo vệ bộ nhớ làm việc của AI.
    \item Bảo vệ ứng dụng trước các sự cố API bằng cơ chế ngắt mạch Circuit Breaker và chính sách thử lại có giãn cách.
\end{itemize}

\subsection{Bài tập thực hành}
\begin{enumerate}[label={\arabic*.}]
    \item \textbf{Refactor State}: Tìm một component React trong dự án cũ đang dùng 3 biến boolean để quản lý dữ liệu mạng và chuyển đổi hoàn toàn sang mô hình \texttt{AsyncState<T>} 4 bước.
    \item \textbf{Phân rã tệp cồng kềnh}: Chọn một tệp mã nguồn dài trên 350 dòng và hướng dẫn AI phân rã thành 3 tệp nhỏ chuyên biệt (types, service, ui) mà không làm gián đoạn ứng dụng.
    \item \textbf{Thiết lập Zod Gatekeeper}: Viết một hàm xác thực dữ liệu đầu vào cho một API endpoint, kiểm tra tối thiểu 4 điều kiện ràng buộc dữ liệu bằng Zod trước khi cho phép lưu vào database.
\end{enumerate}

\end{document}
